\section{导读：从现象级产品到重新理解 AI Native}

本期访谈的入口是妙鸭。2023 年，张月光在阿里做出过妙鸭这样现象级的 AI 写真产品；但两年后他反而说，妙鸭不是 AI Native 产品。这句话是整期最重要的认知转折。它不是否定妙鸭，而是把“用了 AI 技术”和“产品范式被 AI 改写”区分开。

本节先建立全篇的主线：张月光作为互联网时代训练出来的产品经理，经历了大厂、第一次创业、妙鸭、离开阿里、创办沐言智语、探索星眠和 Dokie，最终把问题收束到两个判断。第一，AI Native 产品的核心不是是否用 AI，而是从流程设计转向上下文设计。第二，下一阶段最值得押注的是 Agent，尤其是能够突破个人能力边界的 AI 朋友。

\begin{importantbox}{本期核心命题}
AI Native 不是“没有 AI 做不了”的简单定义。更高阶的定义是：用户输入和模型输出都变得开放，产品经理不能再穷举流程，而要设计上下文、工具、模型组合、反馈和产品心智。
\end{importantbox}

\begin{knowledgebox}{视觉策略说明}
本视频是固定访谈画面，没有 slides、白板或产品演示。按播客笔记标准，正文不重复插入人物帧；封面用于来源识别，正文用流程到上下文、One Way Door、AI 产品组织、Agent 反馈环等概念图来解释内容。
\end{knowledgebox}

\subsection{本章小结}

这期不是“妙鸭复盘”那么简单。它真正讨论的是互联网产品经理如何穿越 AI 范式变化：从流程、按钮和漏斗，走向上下文、模型能力、Agent 行动和用户心智。

